% paper/skeleton_aisec.tex
% ═══════════════════════════════════════════════════════════════════════════
% Kavach — AISec 2026 submission skeleton (ACM CCS workshop, The Hague)
% Deadline: July 24, 2026.  10pp ACM double-column, ANONYMIZED.
%
% Restructured June 10, 2026 per repo audit (KAVACH_SCRUTINY_AISEC_PLAN.md):
%   - Spine: systems/measurement (the hook-gap finding + live interception +
%     representation-mismatch + hybrid-retrieval fix).
%   - §5: the aggregation frontier (veto vs Bayesian vs hybrid under adaptive
%     corruption) — now testable against the REAL Speaker via adaptive_attack v3.
%   - The old "temporal--spatial" section is DEMOTED to a short design
%     principle (§3.5); the implementation does not enforce load-order, so the
%     standalone formal section over-claimed and burned page budget.
%
% Numbers are [TBD] until Parv's Dell runs land (see DELL_BENCHMARK_RUNBOOK.md).
% Replace \anonymise{} contents only in the camera-ready.
% ═══════════════════════════════════════════════════════════════════════════

\documentclass[sigconf,anonymous,review]{acmart}

% AISec is non-archival-friendly; keep rights commands minimal for review.
\settopmatter{printacmref=false}
\renewcommand\footnotetextcopyrightpermission[1]{}
\pagestyle{plain}

\usepackage{booktabs}
\usepackage{amsmath, amssymb}
\usepackage{graphicx}

\newcommand{\anonymise}[1]{\textit{[anon]}}
\newcommand{\TBD}{\textcolor{red}{[TBD]}}
\usepackage{xcolor}

\begin{document}

\title{Kavach: Semantic Pre-Execution Interception for LLM Agent Runtimes}
\subtitle{A Measurement, a Deployment, and the Aggregation Frontier}

\author{\anonymise{Author One}}
\affiliation{\institution{\anonymise{PES University}}\country{\anonymise{India}}}
% NOTE: capstone mentor Prof. Rajesh Banginwar acknowledged in camera-ready only.

\begin{abstract}
We present Kavach, an OpenClaw-native runtime monitor that intercepts every
agent tool call \emph{before execution}, scores its semantic intent against a
curated corpus of attack-pattern descriptions across four specialized
ministers (EXECUTOR / VAULT / CHANNEL / NAVIGATOR), and tracks session-level
trajectory and intent drift. We make three contributions. First, a
\emph{measurement}: we show that OpenClaw's gateway execution path exposes no
pre-execution plugin hook, so no security plugin can intercept there, and we
characterize where interception \emph{is} achievable (the embedded runner).
Second, a \emph{deployment}: live interception on production hardware at
\TBD\,ms p50, with a representation-mismatch analysis explaining why a corpus
authored in tool-call syntax yields \TBD\% FPR on natural-language attacks, and
a hybrid BM25+dense retrieval fix that recovers FPR to \TBD\%. Third, the
\emph{aggregation frontier}: we show that score-based vote aggregation
(including a Bayesian variant) is brittle to an adaptive attacker who corrupts
individual ministers, while a pure-veto aggregator is robust but sacrifices
benign utility; we quantify the FPR/robustness frontier and propose a tuned
hybrid (Bayesian decision path with a veto floor). All artifacts are
reproducible offline; the trajectory monitor and aggregator require no GPU.
\end{abstract}

\maketitle

% ═══════════════════════════════════════════════════════════════════════════
\section{Introduction}
\label{sec:intro}
% - The harness, not the model, is the attack surface (CVE litany).
% - Pre-execution interception is the right control point, but only works
%   where the runtime exposes a hook — our measurement (§4.1) shows this is
%   not free.
% - Positioning vs MXC and ClawGuard (ready-to-paste text in
%   paper/PAPER_NOTES_JUNE3.md §1). ClawGuard's authors explicitly call for
%   "system-level semantic monitoring" — exactly Kavach's layer.
%
% Contributions (revised — three, all backed by repo artifacts):
%   1. Measurement: the OpenClaw gateway pre-execution hook gap (§4.1).
%   2. Deployment: live semantic interception + the representation-mismatch
%      finding + the hybrid-retrieval FPR fix (§4.2–4.4).
%   3. The aggregation frontier: veto vs Bayesian vs hybrid under adaptive
%      minister corruption, measured on real minister-vote dumps (§5).

% ═══════════════════════════════════════════════════════════════════════════
\section{Background}
\label{sec:background}
\subsection{Agent infrastructure attack surface}
% tool calls / MCP / harness shift. Cite Liu, Hou, ETDI, STRIDE.
\subsection{Where guardrails enforce}
% prompt-time / CoT-time / tool-time / reply-time taxonomy.
% Cite LlamaFirewall 2505.03574, AgentSpec, AGrail, ShieldAgent.
\subsection{Closest prior art: ClawGuard and MXC}
% ClawGuard 2604.11790 (rule-based, OpenClaw-native, 0% ASR AgentDojo under
% rules, but no semantics, no session awareness). MXC (OS containment, not a
% semantic monitor). Three complementary layers: MXC → ClawGuard → Kavach.
% Ready-to-paste text in PAPER_NOTES_JUNE3.md §6.

% ═══════════════════════════════════════════════════════════════════════════
\section{Kavach Architecture}
\label{sec:architecture}
\subsection{Overview}
% Pipeline: user msg → seed_intent → action → COMPASS drift → semantic router
% → activated ministers (parallel) → Speaker → BLOCK/ESCALATE/ALLOW.
\subsection{The four ministers}
% Per-minister taxonomy table (MITRE/OWASP), corpus size, threshold.
% Honest note: report per-minister TPR/FPR from the sweep (EXECUTOR strong,
% NAVIGATOR weak on InjecAgent) and pre-empt the "why four?" question with the
% ablation in §5; NAVIGATOR is a trajectory-class detector InjecAgent doesn't
% exercise — show it firing in the native staged-attack test (§4).
\subsection{COMPASS — intent-alignment oracle}
% seed_intent stores the BGE vector of the user message; per-call
% cosine(intent, action); drift threshold. Honest #9 caveat: live seeding.
\subsection{Hybrid retrieval}
% BM25 + dense fused with RRF for SELECTION; confidence = dense cosine ×
% lexical gate (GATE_FLOOR knob). This is the representation-mismatch fix —
% forward-reference §4.3. Cite RRF (Cormack et al.).
\subsection{Speaker and the session monitor}
% The deployed Speaker is a PURE VETO over minister verdicts (any BLOCK ≥
% per-minister threshold → BLOCK), with a trajectory-ceiling case (Case 0)
% that blocks on accumulated session risk independent of any single call.
% Trajectory signals: denial echo (ARM), sequential direction (Trajectory
% Guard), cross-minister escalation (AgentDoG), hard ceiling.
% (Optional 1-paragraph design principle, NOT a full section: pre-execution
% ordering — ministers authoritative before the tool runs. State it as a
% design goal realized by the before_tool_call hook; do not over-formalize.)

% ═══════════════════════════════════════════════════════════════════════════
\section{Deployment and Detection}
\label{sec:eval}
\subsection{The gateway hook gap (measurement)}
\label{sec:hookgap}
% THE headline systems finding. Embedded runner fires before_tool_call;
% gateway/TUI path does not → no plugin can intercept pre-execution there.
% Sentinel-file evidence. Generalize: pre-execution hook should be a mandatory
% MAS primitive. This stands alone even if every benchmark slipped.
\subsection{Live interception (deployment)}
% Hardware (Dell Precision 3660, RTX 4090; Gemma\,4 27B via Ollama; BGE-base
% 768-d). Live BLOCK logged to the tamper-evident ledger, ~820ms.
% Latency: report p50/p95 MEASURED (826/1649 ms) — no "<1s" target as result.
\subsection{Representation mismatch (the instructive negative result)}
\label{sec:repmismatch}
% Native tool-call format: \TBD\% FPR. Natural-language InjecAgent at the
% recall-preserving threshold: \TBD\% FPR. Same recall, very different FPR —
% the corpus is authored in tool-call syntax. This is a lesson-learned the
% workshop values.
\subsection{Hybrid retrieval recovers FPR}
% Re-run with hybrid BM25+dense: FPR \TBD\% at recall \TBD\%. GATE_FLOOR
% ablation table. Tie back to §3.4.
\subsection{Detection results}
% InjecAgent strict/loose recall + FPR (per-minister thresholds + hybrid).
% Native benchmark (clean format). Tamper-evident ledger: /ledger/verify
% detects historical edits (one sentence + appendix).
\subsection{Cost}
% Embeds/call: 1 (embed-once) vs up to 9 naive. Latency percentiles.
% Compare qualitatively to PRISM (~15.8s) and AgentArmor (~20.9s) — NOT
% cross-benchmark ASR ratios.

% ═══════════════════════════════════════════════════════════════════════════
\section{The Aggregation Frontier}
\label{sec:frontier}
% The Pivot-B contribution, now testable against the REAL Speaker.
\subsection{Adaptive attacker model}
% Attacker corrupts K most-foolable ministers per attack action (forces
% ALLOW@high-confidence). Measured on real minister-vote dumps
% (adaptive_attack.py v3, HANDOFF_SCHEMA.md).
\subsection{Veto vs Bayesian vs hybrid}
% The deployed pure-veto Speaker: lowest ASR at every K, but clean FPR above
% the 5% gate. Bayesian: passes FPR, collapses at K=2. Hybrid (Bayesian path +
% veto floor at swept confidence): passes FPR AND tracks veto robustness.
% Frontier table from --sweep-veto. THIS IS THE CONTRIBUTION: we characterize
% the tension, we do not claim a single best aggregator.
\subsection{Minister ablation}
% 1..4 ministers, with/without COMPASS, with/without trajectory. Addresses
% "why four ministers" given NAVIGATOR's weak standalone InjecAgent numbers.

% ═══════════════════════════════════════════════════════════════════════════
\section{Related Work}
\label{sec:related}
% Full comparison table in PAPER_NOTES_JUNE3.md (ClawGuard, AgentArmor,
% LlamaFirewall, MELON, PromptArmor, PRISM, ZEDD, MXC, Kavach). Keep ASR
% numbers same-benchmark only; mark cross-benchmark cells explicitly.

% ═══════════════════════════════════════════════════════════════════════════
\section{Limitations}
\label{sec:limitations}
% - ClawGuard achieves 0% ASR on AgentDojo under rules; we do not match that
%   on AgentDojo (complementary semantic approach; head-to-head is future work
%   pending our AgentDojo run).
% - Tamper-EVIDENT not tamper-PROOF (write-capable attacker; external anchor
%   is future work).
% - Trajectory ceiling needs intent seeding (#9) to reach full risk.
% - Mean-pooling chain leg vs trained Siamese (Trajectory Guard F1 ≤ 0.69).
% - Single LLM backbone (Gemma\,4 27B); cross-model generality is future work.

% ═══════════════════════════════════════════════════════════════════════════
\section{Conclusion}
\label{sec:conclusion}

\bibliographystyle{ACM-Reference-Format}
% \bibliography{kavach}

\end{document}
